% LIGHT v0 system map: energy flow through the Kr/F2 kinetics.
% Every number comes from system_map_numbers.tex, which
% make_system_map.py writes from the model. Build from LIGHT/docs:
%   python3 make_system_map.py && latexmk -pdf system_map.tex
\documentclass[11pt]{article}
\usepackage[a4paper,landscape,margin=0.9cm]{geometry}
\usepackage[T1]{fontenc}
\usepackage[scaled=0.95]{helvet}
\renewcommand{\familydefault}{\sfdefault}
\usepackage{amsmath}
\usepackage{sansmath}
\usepackage{booktabs,tabularx,array}
\usepackage{xcolor}
\usepackage{tikz}
\usetikzlibrary{arrows.meta,calc,positioning}
\pagestyle{empty}
\setlength{\parindent}{0pt}

% LIGHT system-map numbers. GENERATED by make_system_map.py from the
% v0 model (0.1.0). Do not edit. Rerun the script instead.
\providecommand\lmdef[2]{\expandafter\def\csname lm@#1\endcsname{#2}}
\lmdef{e.P}{100.0}
\lmdef{e.pen}{6.4}
\lmdef{e.G1}{55.1}
\lmdef{e.C1}{37.4}
\lmdef{e.C2}{10.4}
\lmdef{e.G6}{3.4}
\lmdef{e.dU}{0.0}
\lmdef{C1.kr}{31.5}
\lmdef{C1.heat}{5.9}
\lmdef{krs.in}{86.6}
\lmdef{krs.G3}{64.5}
\lmdef{krs.G2}{21.7}
\lmdef{krs.left}{0.4}
\lmdef{brs.G3}{74.5}
\lmdef{brs.G2}{25.0}
\lmdef{G3.krf}{40.0}
\lmdef{G3.F}{5.2}
\lmdef{G3.heat}{19.3}
\lmdef{G2.ion}{15.3}
\lmdef{C2.ion}{10.4}
\lmdef{krf.G4}{20.6}
\lmdef{krf.G5}{19.4}
\lmdef{krf.left}{0.0}
\lmdef{brf.G4}{51.6}
\lmdef{brf.G5}{48.4}
\lmdef{krf.eff}{32.5}
\lmdef{G4.hv}{16.8}
\lmdef{G4.F}{2.7}
\lmdef{G4.heat}{1.2}
\lmdef{G5.k2f}{12.9}
\lmdef{G5.heat}{6.5}
\lmdef{k2f.C3}{12.9}
\lmdef{k2f.left}{0.0}
\lmdef{C3.hv}{9.4}
\lmdef{C3.F}{2.5}
\lmdef{C3.heat}{0.9}
\lmdef{out.248}{16.8}
\lmdef{out.415}{9.4}
\lmdef{out.heat}{33.8}
\lmdef{out.elastic}{3.4}
\lmdef{out.ions}{25.7}
\lmdef{out.F}{10.4}
\lmdef{out.other}{0.4}
\lmdef{version}{0.1.0}
\lmdef{p}{1}
\lmdef{xF2}{0.5}
\lmdef{Tg0}{300}
\lmdef{P}{300}
\lmdef{pulse}{1}
\lmdef{tend}{2}
\lmdef{ne}{10^{15}}
\lmdef{N}{2.45{\times}10^{19}}
\lmdef{tau}{6.99}
\lmdef{Edep}{0.30}
\lmdef{Te}{1.31}
\lmdef{TgEnd}{519}
\lmdef{pEnd}{1.74}
\lmdef{F2burn}{830}
\lmdef{F2left}{0.3}
\lmdef{krsFormed}{1.64{\times}10^{17}}
\lmdef{krfFormed}{1.22{\times}10^{17}}
\lmdef{notes.krs}{92.5}
\lmdef{notes.krp}{25.7}
\lmdef{notes.emit}{26.2}
\lmdef{om.quench}{10.9}
\lmdef{om.kr2}{17.1}
\lmdef{om.ions}{17.6}
\lmdef{om.electron}{15.4}
\lmdef{om.bc}{10.1}
\lmdef{om.f2e}{20.0}
\lmdef{om.all}{5.9}
\lmdef{om.v0}{16.8}

\makeatletter
\newcommand\lm[1]{\ifcsname lm@#1\endcsname\csname lm@#1\endcsname
  \else\lm@missing{#1}\fi}
\newcommand\lm@missing[1]{\errmessage{system_map_numbers.tex has no
  number '#1'. Rerun make_system_map.py}}
\makeatother

% categorical colors, checked with the dataviz validator (all pairs,
% white surface): blue, amber, aqua, violet, red. Ions and F atoms are
% neutral greys (stored energy). Every band and tag also carries text.
\definecolor{cE}{HTML}{2A78D6}   % electron gas
\definecolor{cX}{HTML}{EDA100}   % Kr* excitation
\definecolor{cB}{HTML}{1BAF7A}   % KrF(B) and 248 nm
\definecolor{cT}{HTML}{4A3AA7}   % Kr2F* and 415 nm
\definecolor{cH}{HTML}{E34948}   % heat
\definecolor{cI}{HTML}{52514E}   % ions (stored)
\definecolor{cF}{HTML}{A8A69E}   % F atoms (stored)
\definecolor{ink}{HTML}{0B0B0B}
\definecolor{ink2}{HTML}{52514E}
\definecolor{muted}{HTML}{898781}

% chemistry and units in running text or math
\newcommand{\unit}[1]{\ensuremath{\,\mathrm{#1}}}
\newcommand{\el}{\ensuremath{\mathrm{e}^-}}
\newcommand{\Krs}{Kr\ensuremath{^*}}
\newcommand{\Krp}{Kr\ensuremath{^+}}
\newcommand{\Krtp}{Kr\ensuremath{_2^+}}
\newcommand{\Krts}{Kr\ensuremath{_2^*}}
\newcommand{\KtF}{Kr\ensuremath{_2}F\ensuremath{^*}}
\newcommand{\Ft}{F\ensuremath{_2}}
\newcommand{\ra}{\ensuremath{\to}\ }
\newcommand{\hnu}{\ensuremath{h\nu}}
\newcommand{\Edep}{\ensuremath{E_\mathrm{dep}}}
\newcommand{\pc}[1]{\lm{#1}\,\%}

\tikzset{
  band/.style={line width={max(0.6pt,(#1)*0.03mm)},
    -{Triangle[width=2.2pt 1.8,length=2pt 1.0]}},
  thin arrow/.style={line width=0.6pt,
    -{Triangle[width=2.2pt 1.8,length=2pt 1.0]}},
  blab/.style={font=\scriptsize, text=ink, inner sep=0.8pt, fill=white,
    fill opacity=0.85, text opacity=1},
  nohyph/.style={execute at begin node={\hyphenpenalty=10000
    \exhyphenpenalty=10000}},
  proc/.style={nohyph, draw=ink2!55, fill=white, rounded corners=2pt,
    text width=2.55cm, align=left, font=\scriptsize, inner sep=3pt,
    text=ink},
  spec/.style={nohyph, draw=#1, fill=#1!12, rounded corners=5pt,
    align=center, font=\scriptsize, text=ink, inner sep=3pt,
    text width=2.1cm, line width=0.8pt},
  inbox/.style={nohyph, draw=ink2!55, fill=white, rounded corners=2pt,
    text width=2.7cm, align=left, font=\scriptsize, inner sep=4pt,
    text=ink},
  tag/.style={draw=#1, fill=#1!14, rounded corners=1.5pt,
    font=\tiny, text=ink, inner sep=1.6pt},
  hdr/.style={font=\footnotesize\bfseries, text=ink2},
  note/.style={nohyph, font=\tiny, text=ink2},
}

\begin{document}
\sansmath

{\Large\bfseries LIGHT v0 system map: energy flow through the
Kr/F$_2$ kinetics}\par
\vspace{2pt}
{\small Reference point: \lm{p}\unit{atm} of Kr with \pc{xF2} F$_2$ at
\lm{Tg0}\unit{K}, $P=\lm{P}\unit{kW\,cm^{-3}}$ for \lm{pulse}\unit{\mu s}
so that $\Edep=\lm{Edep}\unit{J\,cm^{-3}}$, $n_e=\lm{ne}\unit{cm^{-3}}$
held, run to \lm{tend}\unit{\mu s}. Band width is proportional to the
energy flow. Numbers are percent of \Edep, rounded to 0.1.}\par
\vspace{2pt}

\begin{tikzpicture}[x=1cm,y=1cm]
% ================================================================ frames
\node[hdr] at (1.55,15.5) {INPUT};
\node[hdr] at (26.0,15.5) {OUTPUT at $t=\lm{tend}\unit{\mu s}$};
\node[note] at (26.0,15.12) {percent of \Edep, sums to 100};
\draw[draw=ink2!35, fill=black!2.5, rounded corners=6pt]
  (3.5,0.25) rectangle (23.9,15.25);
\node[hdr, anchor=west] at (3.7,14.88)
  {UNIFORM INITIAL GAS\quad 0-D, constant volume};

% ================================================================ inputs
\node[inbox, line width=0.9pt, draw=cE] (pow) at (1.55,12.4) {%
  \textbf{Power in}\\
  $P=\lm{P}\unit{kW\,cm^{-3}}$\\ for \lm{pulse}\unit{\mu s}, uniform,\\
  into the electrons\\
  $\Edep=\lm{Edep}\unit{J\,cm^{-3}}$\\
  \textcolor{ink2}{range 100 to 500}};
\node[inbox, dashed] (em) at (1.55,9.6) {%
  \textbf{or an EM wave}\\
  $E_\mathrm{rms}$ at frequency $f$\\
  absorbed by electrons\\
  \textcolor{ink2}{not used at this point}};
\node[inbox] (inj) at (1.55,7.0) {%
  \textbf{$n_e$ injection}\\
  $n_e=\lm{ne}\unit{cm^{-3}}$, held\\
  \textcolor{ink2}{range $10^{14}$ to $10^{16}$}};
\node[inbox] (fill) at (1.55,4.2) {%
  \textbf{Initial gas fill}\\
  Kr $+$ \pc{xF2} F$_2$\\ \lm{p}\unit{atm}, \lm{Tg0}\unit{K}\\
  \mbox{$N=\lm{N}\unit{cm^{-3}}$}\\
  \textcolor{ink2}{range 0.1 to 2\unit{atm}}};

% ================================================================ electron gas
\node[draw=cE, fill=cE!10, rounded corners=5pt, line width=0.9pt,
  minimum width=1.7cm, minimum height=10.1cm, anchor=north]
  (eg) at (4.65,14.4) {};
\node[align=center, font=\scriptsize, text=ink, text width=1.6cm,
  anchor=north] at ([yshift=-4pt]eg.north) {\textbf{electron}\\
  \textbf{gas}\\[3pt] Maxwellian\\[3pt] $T_e=\lm{Te}\unit{eV}$\\
  \textcolor{ink2}{pulse mean}\\[3pt] $U_e=\tfrac32 n_e T_e$};
\node[align=center, note, text width=1.6cm, anchor=south]
  at ([yshift=3pt]eg.south) {in: power \lm{e.P}\\ and Penning
  \lm{e.pen}};

% ================================================================ e-impact rows
\node[note, anchor=south west] at (6.2,14.15)
  {$Q=k\,n_e[\mathrm{Kr}]$, $k=\langle\sigma v\rangle$ over a Maxwellian,
  Biagi v7.1 cross sections};
\node[proc] (G1) at (7.6,13.3) {%
  \textbf{G1}\quad \el{} + Kr\\ \ra \el{} + \Krs(5s)\\
  \textcolor{ink2}{threshold 9.915\unit{eV}}};
\node[proc] (C1) at (7.6,11.05) {%
  \textbf{C1}\quad \el{} + Kr\\ \ra Kr(5p, 4d, \dots)\\
  \textcolor{ink2}{cascades to \Krs(5s)}};
\node[proc] (C2) at (7.6,6.3) {%
  \textbf{C2}\quad \el{} + Kr\\ \ra \Krp{} + 2\el\\
  \textcolor{ink2}{threshold 13.996\unit{eV}}};
\node[proc] (G6) at (7.6,4.75) {%
  \textbf{G6}\quad elastic\\ \el{} + Kr\\
  \textcolor{ink2}{$(2m/M)(1-T_g/T_e)$}};

% ================================================================ Kr*
\node[spec=cX] (KS) at (11.05,12.2) {%
  \textbf{\Krs(5s)}\\ formed\\ $\lm{krsFormed}\unit{cm^{-3}}$};

% ================================================================ G3 G2
\node[proc] (G3) at (14.45,13.3) {%
  \textbf{G3}\quad \Krs{} + \Ft\\ \ra KrF(B) + F\\
  \textcolor{ink2}{$8.1{\times}10^{-10}\unit{cm^3s^{-1}}$}\\
  takes \pc{brs.G3} of \Krs};
\node[proc] (G2) at (14.45,10.4) {%
  \textbf{G2}\quad \Krs{} + \Krs\\ \ra \Krp{} + Kr + \el\\
  \textcolor{ink2}{$5.0{\times}10^{-10}\unit{cm^3s^{-1}}$}\\
  takes \pc{brs.G2} of \Krs};

% ================================================================ KrF(B), Kr+
\node[spec=cB] (KF) at (17.85,13.3) {%
  \textbf{KrF(B)}\\ formed\\ $\lm{krfFormed}\unit{cm^{-3}}$};
\node[note, align=center, anchor=north] at ([yshift=-2pt]KF.south)
  {\pc{krf.eff} of \Edep\\ in 5.00-eV quanta};
\node[spec=cI] (KI) at (17.85,6.3) {%
  \textbf{\Krp{} ions}\\ {\tiny stored, not recombined in v0}};

% ================================================================ G4 G5 Kr2F* C3
\node[proc, text width=2.6cm] (G4) at (21.35,14.3) {%
  \textbf{G4}\quad KrF(B)\\ \ra Kr + F + \hnu\\
  \textcolor{ink2}{248\unit{nm}, $\tau=\lm{tau}\unit{ns}$}\\
  \pc{brf.G4} of KrF(B)};
\node[proc, text width=2.6cm] (G5) at (21.35,11.9) {%
  \textbf{G5}\quad KrF(B) + 2Kr\\ \ra \KtF{} + Kr\\
  \textcolor{ink2}{$6.0{\times}10^{-31}\unit{cm^6s^{-1}}$}\\
  \textcolor{ink2}{$\times\,(T_g/300)^{-3.32}$}\\
  \pc{brf.G5} of KrF(B)};
\node[spec=cT] (K2) at (21.35,9.8) {%
  \textbf{\KtF}\\ {\unsansmath$4^2\Gamma$} trimer};
\node[proc, text width=2.6cm] (C3) at (21.35,7.95) {%
  \textbf{C3}\quad \KtF\\ \ra 2Kr + F + \hnu\\
  \textcolor{ink2}{near 415\unit{nm}}\\
  \textcolor{ink2}{$\tau=220\unit{ns}$}};

% ================================================================ bands
% inputs into the gas
\draw[band=\lm{e.P}, cE] (pow.east) -- (pow.east -| eg.west)
  node[blab, midway, above=2.2pt] {\lm{e.P}};
\draw[thin arrow, cE, dashed] (em.east) -- (em.east -| eg.west);
\draw[thin arrow, ink2] (inj.east) -- (inj.east -| eg.west)
  node[blab, midway, above=1pt, font=\tiny] {sets $n_e$};
\draw[thin arrow, ink2] (fill.east) -- (fill.east -| 3.5,0)
  node[blab, midway, above=1pt, font=\tiny] {fills};
% electron gas out to the four electron-impact rows
\foreach \n/\k in {G1/e.G1, C1/e.C1, C2/e.C2, G6/e.G6}{
  \draw[band=\lm{\k}, cE] (\n.west -| eg.east) -- (\n.west)
    node[blab, midway, above=1.5pt] {\lm{\k}};
}
% into Kr*
\draw[band=\lm{e.G1}, cX] (G1.east) -- (KS.north west)
  node[blab, midway, above right=0pt and 0pt] {\lm{e.G1}};
\draw[band=\lm{C1.kr}, cX] (C1.east) -- (KS.south west)
  node[blab, midway, below right=0pt and 0pt] {\lm{C1.kr}};
% out of Kr*
\draw[band=\lm{krs.G3}, cX] (KS.north east) -- (G3.west)
  node[blab, midway, above left=0pt and 0pt] {\lm{krs.G3}};
\draw[band=\lm{krs.G2}, cX] (KS.south east) -- (G2.west)
  node[blab, midway, below left=0pt and 0pt] {\lm{krs.G2}};
% G3 into KrF(B)
\draw[band=\lm{G3.krf}, cB] (G3.east) -- (KF.west)
  node[blab, midway, above=2pt] {\lm{G3.krf}};
% G2 into the ions and back to the electrons
\draw[band=\lm{G2.ion}, cI] (G2.south east) -- (KI.north west)
  node[blab, midway, right=3pt] {\lm{G2.ion}};
\draw[band=\lm{e.pen}, cE, rounded corners=5pt]
  (G2.south west) -- (G2.south west |- 0,9.15) -- (eg.east |- 0,9.15)
  node[blab, pos=0.5, below=2pt] {Penning returns \lm{e.pen} to the
  electrons, 5.83\unit{eV} per event};
\draw[band=\lm{C2.ion}, cI] (C2.east) -- (KI.west)
  node[blab, pos=0.5, above=1.5pt] {\lm{C2.ion}};
% out of KrF(B)
\draw[band=\lm{krf.G4}, cB] (KF.north east) -- (G4.west)
  node[blab, midway, above left=0pt and 0pt] {\lm{krf.G4}};
\draw[band=\lm{krf.G5}, cB] (KF.south east) -- (G5.west)
  node[blab, midway, below left=0pt and 0pt] {\lm{krf.G5}};
% light, trimer and ions to the outputs
\draw[band=\lm{G4.hv}, cB] (G4.east) -- (G4.east -| 24.35,0)
  node[blab, midway, above=2pt] {\lm{G4.hv}};
\draw[band=\lm{G5.k2f}, cT] (G5.south) -- (K2.north)
  node[blab, midway, left=3pt] {\lm{G5.k2f}};
\draw[band=\lm{k2f.C3}, cT] (K2.south) -- (C3.north)
  node[blab, midway, left=3pt] {\lm{k2f.C3}};
\draw[band=\lm{C3.hv}, cT] (C3.east) -- (C3.east -| 24.35,0)
  node[blab, midway, above=1.5pt] {\lm{C3.hv}};
\draw[band=\lm{out.ions}, cI] (KI.east) -- (KI.east -| 24.35,0)
  node[blab, midway, above=2.2pt] {\lm{out.ions}};

% ================================================================ loss tags
\node[tag=cH, anchor=north east] at ([yshift=-2pt]C1.south east)
  {heat \lm{C1.heat}};
\node[tag=cH, anchor=north east] at ([yshift=-2pt]G6.south east)
  {heat \lm{e.G6}, elastic};
\node[tag=cF, anchor=north west] at ([yshift=-2pt]G3.south west)
  {F atoms \lm{G3.F}};
\node[tag=cH, anchor=north east] at ([yshift=-2pt]G3.south east)
  {heat \lm{G3.heat}};
\node[tag=cF, anchor=north west] at ([yshift=-2pt]G4.south west)
  {F atoms \lm{G4.F}};
\node[tag=cH, anchor=north east] at ([yshift=-2pt]G4.south east)
  {heat \lm{G4.heat}};
\node[tag=cH, anchor=north east] at ([yshift=-2pt]G5.south east)
  {heat \lm{G5.heat}};
\node[tag=cF, anchor=north west] at ([yshift=-2pt]C3.south west)
  {F atoms \lm{C3.F}};
\node[tag=cH, anchor=north east] at ([yshift=-2pt]C3.south east)
  {heat \lm{C3.heat}};

% ================================================================ not in v0
\node[draw=ink2!60, dashed, rounded corners=3pt, fill=white,
  align=left, font=\scriptsize, text=ink, anchor=south west,
  text width=10.0cm, inner sep=5pt] at (5.95,0.5) {%
  \textbf{Kr/F$_2$ channels not in v0.} The 248-nm share of \Edep{}
  with one optional group switched on. Measured with
  \texttt{run\_v0.py -{}-omissions}.\\[3pt]
  \begin{tabular}{@{}l@{\hspace{3pt}}r@{\hspace{16pt}}l@{\hspace{3pt}}r
    @{\hspace{16pt}}l@{\hspace{3pt}}r@{}}
  v0 as built & \pc{om.v0} & KrF B/C mixing & \pc{om.bc} &
    ion recombination & \pc{om.ions}\\
  collisional quench & \pc{om.quench} & \Krts{} excimer & \pc{om.kr2} &
    \el{} on \Krs & \pc{om.electron}\\
  \textbf{all groups on} & \textbf{\pc{om.all}} & & &
    \el{} on \Ft & \pc{om.f2e}\\
  \end{tabular}};

% ================================================================ legend
\begin{scope}[shift={(16.7,0.55)}, font=\scriptsize, text=ink]
  \node[anchor=west, font=\scriptsize\bfseries, text=ink2] at (0,3.25)
    {Band color is where the energy sits};
  \foreach \c/\t [count=\i] in {cE/electron gas, cX/\Krs{} excitation,
      cB/KrF(B) and 248-nm light, cT/\KtF{} and 415-nm light,
      cI/\Krp{} ions}{
    \draw[\c, line width=1.3mm] (0,3.25-0.5*\i) -- ++(0.55,0);
    \node[anchor=west] at (0.62,3.25-0.5*\i) {\t};
  }
  \node[tag=cH, anchor=west] at (4.15,2.75) {heat};
  \node[anchor=west] at (4.85,2.75) {turns into heat};
  \node[tag=cF, anchor=west] at (4.15,2.25) {F atoms};
  \node[anchor=west] at (5.3,2.25) {left in free F};
  \draw[ink2, line width=0.6pt, dashed] (4.15,1.75) -- ++(0.55,0);
  \node[anchor=west] at (4.75,1.75) {not active here};
  \draw[ink2, line width={10*0.03mm}] (4.15,1.25) -- ++(0.55,0);
  \node[anchor=west] at (4.75,1.25) {10\,\% band};
  \draw[ink2, line width={50*0.03mm}] (4.15,0.75) -- ++(0.55,0);
  \node[anchor=west] at (4.75,0.75) {50\,\% band};
\end{scope}

% ================================================================ ledger
\begin{scope}[font=\scriptsize, text=ink]
  \newcommand{\ledger}[5]{% y, color, label, key, bar options
    \node[anchor=south west, inner sep=1pt] at (24.4,#1) {#3};
    \node[anchor=south east, inner sep=1pt, font=\scriptsize\bfseries]
      at (27.7,#1) {\lm{#4}};
    \draw[ink2!30] (24.45,#1-0.1) -- (27.65,#1-0.1);
    \fill[#2, #5] (24.45,#1-0.1) rectangle ++({3.2*(\lm{#4})/100},-0.14);
  }
  \ledger{14.3}{cB}{248-nm emission}{out.248}{}
  \ledger{7.95}{cT}{415-nm emission}{out.415}{}
  \ledger{6.3}{cI}{stored in \Krp{} ions}{out.ions}{}
  \ledger{5.0}{cH}{reaction heat}{out.heat}{}
  \ledger{4.0}{cH}{elastic losses}{out.elastic}{opacity=0.55}
  \ledger{3.0}{cF}{stored in F atoms}{out.F}{}
  \ledger{2.0}{muted}{other stores}{out.other}{}
  \node[note, anchor=north west, text width=3.25cm, align=left]
    at (24.4,1.45) {The heat and F tags inside the gas add up to these
    totals. Ledger closure is below $10^{-12}$ of \Edep.};
\end{scope}
\end{tikzpicture}

\vspace{1pt}
{\small The notes' outputs as flows, which overlap: to \Krs{}
\pc{notes.krs}, to \Krp{} \pc{notes.krp}, to KrF(B) \pc{krf.eff}, heat
\pc{out.heat}, elastic losses \pc{out.elastic}, emission \pc{notes.emit}.
At the end of the run $T_g=\lm{TgEnd}\unit{K}$ and $p=\lm{pEnd}\unit{atm}$.
F$_2$ falls below 10\,\% of the fill at \lm{F2burn}\unit{ns}.}\par
{\scriptsize\color{ink2} Generated from LIGHT v0 (\lm{version}) by
\texttt{docs/make\_system\_map.py}. No number on this page is typed by
hand. Key and sources on page 2.}

\newpage
% ===================================================================== page 2
{\Large\bfseries Key to the map}\par
\vspace{6pt}
\begin{minipage}[t]{0.44\linewidth}
\small
\textbf{How to read it.} Each band carries energy. Its width is
proportional to that energy, and its number is the energy as a percent
of the deposited energy \Edep. A band's color is the store the energy
sits in while it travels. Red tags are energy that turns into heat at
that step. Grey tags are energy left behind in free fluorine atoms,
because every harpoon event breaks an F$_2$ bond. Dashed outlines mark
what is not active at the reference point.\par
\medskip
Energies are counted at the formation energies below, measured from
ground-state Kr and F$_2$ at rest. What enters a box therefore equals
what leaves it. The script checks that for every box before it writes
a number. Printed values are rounded to 0.1, so a box can be off by 0.1
between its inputs and its outputs.\par
\medskip
The electrons get their energy from the power and from Penning
ionization, which hands 5.83\unit{eV} per event back to them. They lose
it to excitation, ionization and elastic collisions. Energy can
therefore pass through the electrons more than once, and the notes'
flows can add up to more than 100\,\%. Only the output ledger is a
partition. Its last line, other stores, is what is still held in \Krs,
KrF(B), \KtF{} and the electron gas when the run stops.\par
\medskip
\textbf{Formation energies} in eV above ground-state Kr and F$_2$\par
\smallskip
\begin{tabular}{@{}lllll@{}}
\toprule
\Krs(5s) & \Krp & KrF(B) & \KtF & F\\
\midrule
9.915 & 13.996 & 6.15 & 4.10 & 0.80\\
\bottomrule
\end{tabular}\par
\smallskip
The photons carry 5.00\unit{eV} at 248\unit{nm} and 3.0\unit{eV} near
415\unit{nm}.\par
\medskip
\textbf{State at the reference point}\par
\smallskip
\begin{tabular}{@{}ll@{}}
\toprule
$T_e$, pulse mean & \lm{Te}\unit{eV}\\
$T_g$ and $p$ at \lm{tend}\unit{\mu s} & \lm{TgEnd}\unit{K}, \lm{pEnd}\unit{atm}\\
\Krs{} formed & $\lm{krsFormed}\unit{cm^{-3}}$\\
KrF(B) formed & $\lm{krfFormed}\unit{cm^{-3}}$\\
F$_2$ below 10\,\% of the fill & at \lm{F2burn}\unit{ns}\\
F$_2$ left at \lm{tend}\unit{\mu s} & \pc{F2left}\\
\bottomrule
\end{tabular}\par
\medskip
\textbf{Regenerate} after any change to the model:\par
\smallskip
{\footnotesize\ttfamily cd LIGHT/docs\\
python3 make\_system\_map.py\\
latexmk -pdf system\_map.tex}
\end{minipage}\hfill
\begin{minipage}[t]{0.53\linewidth}
\small
\textbf{The v0 reactions.} G-rows are the six in
\texttt{KrF\_model\_notes\_with\_Gary.rtf}. C-rows are the closures that
the notes' outputs need.\par
\smallskip
\renewcommand{\arraystretch}{1.15}
\begin{tabularx}{\linewidth}{@{}l>{\raggedright\arraybackslash}p{3.7cm}
  >{\raggedright\arraybackslash}p{3.3cm}>{\raggedright\arraybackslash}X@{}}
\toprule
& reaction & rate & source\\
\midrule
G1 & \el{} + Kr \ra \el{} + \Krs(5s) & Maxwellian $\langle\sigma v\rangle$ &
  Biagi v7.1, the whole 5s manifold at 9.915\unit{eV}\\
G2 & \Krs{} + \Krs{} \ra \Krp{} + Kr + \el &
  $5.0\times10^{-10}\unit{cm^3s^{-1}}$ & Kannari 1985, Penning products\\
G3 & \Krs{} + \Ft{} \ra KrF(B) + F & $8.1\times10^{-10}\unit{cm^3s^{-1}}$
  & Kannari 1985 Table I\\
G4 & KrF(B) \ra Kr + F + \hnu, 248\unit{nm} & $\tau=\lm{tau}\unit{ns}$ &
  Kannari 1985. The notes say about 4\,ns\\
G5 & KrF(B) + Kr + Kr \ra \KtF{} + Kr &
  $6.0\times10^{-31}\unit{cm^6s^{-1}}\times(T_g/300)^{-3.32}$ &
  Xu, Gadomski and Setser 1993, with Shui's $T_g$ law\\
G6 & \el{} + Kr elastic & power $\propto(2m/M)(1-T_g/T_e)$ &
  Biagi v7.1 momentum transfer\\
C1 & \el{} + Kr \ra Kr(5p, 4d, higher), then \Krs & Maxwellian
  $\langle\sigma v\rangle$ & Biagi v7.1, prompt cascade to 5s\\
C2 & \el{} + Kr \ra \Krp{} + 2\el & Maxwellian $\langle\sigma v\rangle$ &
  Biagi v7.1, 13.996\unit{eV}\\
C3 & \KtF{} \ra 2Kr + F + \hnu & $\tau=220\unit{ns}$ & Xu, Gadomski and
  Setser 1993\\
\bottomrule
\end{tabularx}\par
\medskip
\textbf{Not in v0.} The optional groups carry the rest of the pure
Kr/F$_2$ chemistry in SPECK/REBUILD. The last column is the 248-nm
share of \Edep{} at the reference point with that group switched on. It
is \pc{om.v0} in v0 as built.\par
\smallskip
\begin{tabularx}{\linewidth}{@{}l>{\raggedright\arraybackslash}X r@{}}
\toprule
group & what it adds & 248 nm\\
\midrule
quench & KrF(B) and \KtF{} quenched by Kr, F$_2$ and electrons &
  \pc{om.quench}\\
bc & the KrF C state, which mixes fast with B and radiates ten times
  slower & \pc{om.bc}\\
kr2 & \Krs{} + 2Kr \ra \Krts, its 147-nm light and its own harpoon &
  \pc{om.kr2}\\
ions & \Krp{} \ra \Krtp{} and recombination back to \Krs &
  \pc{om.ions}\\
electron & stepwise ionization and superelastic loss of \Krs &
  \pc{om.electron}\\
f2e & electrons on F$_2$, meaning attachment, dissociation and
  vibration, and ion-ion recombination to KrF(B) & \pc{om.f2e}\\
all & every group at once & \pc{om.all}\\
\bottomrule
\end{tabularx}
\end{minipage}

\end{document}
